\section{Agent 市场分层：personal、role-based、customer-facing}

讲者把 Agent 生态划分为三类：personal agents（如 ChatGPT/Gemini）、role-based agents（如 coding/legal assistants）、customer-facing agents（企业对外服务 Agent）。这个框架的价值在于，三类 Agent 的成功函数不同，不能用同一套评测和产品方法。

\begin{table}[H]
\centering
\caption{三类 Agent 的目标函数差异}
\begin{tabularx}{\textwidth}{>{\raggedright\arraybackslash}p{2.8cm} >{\raggedright\arraybackslash}p{3.2cm} >{\raggedright\arraybackslash}X >{\raggedright\arraybackslash}p{3.2cm}}
\toprule
类型 & 主要用户 & 主任务 & 首要评价指标 \\
\midrule
personal agents & 个人终端用户 & 信息获取、写作、探索 & 可用性与泛化能力 \\
role-based agents & 专业岗位人员 & 代码、法务、研究支持 & 专业效率与正确率 \\
customer-facing agents & 企业外部客户 & 服务、故障排查、交易变更 & 可验证解决率、升级率、合规率 \\
\bottomrule
\end{tabularx}
\end{table}

\begin{knowledgebox}{为什么课程里要单独讲 customer-facing agent}
这类 Agent 面向的是“外部真实客户 + 企业系统 of record + 高并发会话”。其难点不仅是模型能力，更是跨渠道状态一致性、合规动作约束、SLA 与回退策略设计。
\end{knowledgebox}

\begin{importantbox}{讲者判断}
\textit{``Agents are to AI as the website was to the internet and apps were to mobile.''}
对应到企业世界，这句话的含义是：未来每个企业都会有自己的对外 Agent，正如每个企业今天都有官网与 App。
\end{importantbox}

\subsection{本章小结}

“Agent 三分法”帮助我们避免方法误配。课程后续所有工程细节，都是为了让 customer-facing agent 满足企业级外部服务要求，而不是泛化聊天场景。

